\section{Experimental Protocol}
\label{sec:experiments}
\draftnote{This section specifies intended experiments. Dataset sizes and
hyperparameters below are pilot configurations, not completed runs or
validated choices.}

\subsection{Models and disjoint development splits}
The primary model is Qwen2.5-VL-7B-Instruct~\cite{qwen25vl}. We will test
transfer to Qwen3-VL-8B-Instruct~\cite{qwen3vl} and Gemma 4 E4B-it~\cite{gemma4}.
The two Qwen models test transfer across generations, while Gemma 4 tests
a different model family. E4B has 4.5B effective parameters and 8B total
including embeddings; absolute scores across models are not capacity-matched.
We will pin exact model and
processor revisions and use the same image preprocessing, prompts, decoding
settings, and evaluation examples for every method on a given model.
Current configurations use batch size one, BF16, eager attention, a
448-square input preprocessing step, and greedy cached generation.
The model-specific processor may apply further resizing; the effective
image and token dimensions must be recorded.

COCO train images~\cite{coco} supply disjoint discovery, fitting, and
hyperparameter-tuning groups. An initial development allocation is 200
original images per split, with at most 100 discovery records processed
by the starting configuration. All questions and edited variants of an
original image remain in the same group. COCO validation images and
external evaluations are checked for original-image overlap rather than
assumed independent because benchmark names differ.

Initial \method settings are $K=8$, $\lambda_c=1$, $\eta=1$, and
$\epsilon=0.05$. Candidate decoder layers are $\{8,12,16,20,24\}$ for the
two Qwen models and $\{5,11,23,35,41\}$ for Gemma 4 E4B; the latter five
are full-attention layers. All query heads in these layers are considered.
Gains use Adam for two balanced epochs
with learning rate $0.03$, $\rho=0.25$, $\lambda_{\mathrm{KL}}=1$, and
$\lambda_2=0.01$. These settings are starting points. Any tuning will use
development data and a declared budget before final evaluation.

\subsection{Evidence and context diagnostics}
The primary diagnostic follows Table~\ref{tab:conditions}. Each original
image group should support target-present and target-absent conditions,
crossed with the presence or absence of an associated contextual cue.
Target identity, visibility, and position should be preserved when changing
context; context should be preserved as far as possible when removing the
target. Image edits require quality checks for residual objects, boundary
artifacts, unintended semantic changes, and ambiguous labels.

This semantic-context diagnostic is a \emph{planned addition}; the current
gray-occlusion builder does not supply it. We will also evaluate the
unchanged Causal-HalBench test material and its official metrics rather
than repurpose that benchmark as method-selection data~\cite{causalhalbench}.
Native benchmark fields and metric names will be preserved.

We will report accuracy for each condition and the fraction of complete
four-condition groups for which every answer is correct. For controlled
target changes, both answers must be correct; merely flipping the answer
does not count as success. For context-only changes, both answers must
again be correct; constant wrong answers receive no credit.

\paragraph{Separating discrimination from response bias.}
For forced-choice answers, we will compare false-positive rates at
development-calibrated operating points with similar true-positive rates,
and report the achieved test rates rather than assume exact matching.
Thresholds are chosen only on the tuning split. ROC curves provide a
complementary threshold-sweep diagnostic. A base-model margin-offset
baseline changes the Yes/No threshold without changing representations.
This control is essential: a global score shift must not be mistaken for
improved evidence discrimination.

\subsection{Natural images, information retention, and capability}
POPE random, popular, and adversarial subsets provide natural-image
object-existence evaluation~\cite{pope}. We report accuracy, precision,
recall, F1, affirmative-answer rate, and invalid-response rate. Generated
answers use the same parser; malformed or noncommittal answers are not
silently discarded.

On an intended 2,000-image COCO validation subset, we will measure
CHAIR$_s$ and CHAIR$_i$ alongside correct object recall, distinct correct
object mentions, and actual caption length~\cite{chair,coco}.
Generation limits are matched, and multiple steering strengths are used to
inspect the hallucination--recall relationship. The CHAIR object taxonomy
does not establish correctness of every attribute or relation.
MME supplies an additional capability-retention check~\cite{mme}.
MMHal-Bench supplies 96 free-response hallucination questions~\cite{mmhal}.
We will use its official response format and one fixed judge version and
scoring prompt for all methods on a model. Judge outputs and response
lengths will be retained; its score is not treated as a Yes/No accuracy.

\subsection{Baselines and implementation parity}
The comparison includes the unmodified model, threshold calibration,
VTI, GrAInS, REVIS, and RUDDER. HONES and the steering method studied in
Seeing or Knowing? are close mechanistic comparators; PURGE and
object-focused preference optimization provide relevant training-based
comparisons where computational and data budgets permit.

The current four steering ports use a common runtime but contain material
adaptations. GrAInS attributes fused decoder-input embeddings from a zero
baseline; VTI uses a fixed visual grid; REVIS uses a fixed neutral-response
choice; and RUDDER exposes an explicit first-token intervention choice.
Common COCO Yes/No preference pairs also differ from original caption
preference data. Therefore, the present ports must be labeled
\emph{adapted}, and checked against author implementations on shared small
inputs before results are presented as reproductions.

Every method receives a declared development and tuning budget. We account
for \method's segmentation masks, paired edits, labels, and discovery
computation, not only the number of fitted parameters. Published scores
from different experimental settings will not be inserted into the same
comparison table.

\subsection{Ablations and uncertainty}
Required ablations replace causal selection with random or activation-RMS
selection while keeping the same $K$, candidate pool, fitting data, and
optimization budget. Further tests remove the control penalty or the KL
term, vary $K$ and the gain bound, and compare last-token with all-prefill
intervention. A residual-scaling control tests whether a generic magnitude
change explains the effect.

We will repeat discovery and fitting with at least three specified seeds
and report all runs. Paired confidence intervals resample original image
groups, keeping correlated questions and variants together; seed
variability is reported separately. Inference timing uses synchronized GPU
measurements after warm-up and is accompanied by peak memory, output
length, and the actual device configuration.
